%----------------------------------------------------------------------------------------
%   Statement of Purpose - Nabil Ahmad (European Universities & Research Groups)
%   Tailored for European MSc & Research Groups (TUM, ETH, EPFL, KTH, etc.)
%----------------------------------------------------------------------------------------
\documentclass[10pt,a4paper]{article}

\usepackage[utf8]{inputenc}
\usepackage[margin=1.5cm,top=1.3cm,bottom=1.3cm]{geometry}
\usepackage{titlesec}
\usepackage{hyperref}
\usepackage{xcolor}

\hypersetup{
    colorlinks=true,
    linkcolor=blue!70!black,
    urlcolor=blue!70!black
}

\titleformat{\section}{\normalsize\bfseries\uppercase}{}{0em}{}[\titlerule]
\titlespacing*{\section}{0pt}{6pt}{3pt}

\setlength{\parindent}{15pt}
\setlength{\parskip}{3pt}

\begin{document}
\pagestyle{empty}

%----------------------------------------------------------------------------------------
%   HEADER
%----------------------------------------------------------------------------------------
\begin{center}
    {\large \textbf{STATEMENT OF PURPOSE}} \\[2pt]
    \textbf{Candidate for Master’s (MSc) $\cdot$ Software Engineering, Human-Computer Interaction \& Applied AI} \\[2pt]
    \small \textbf{Nabil Ahmad} \quad$|$\quad Peshawar, Pakistan \quad$|$\quad +92 304 1237882 \quad$|$\quad \href{mailto:nabilroghani96@gmail.com}{nabilroghani96@gmail.com} \\[1pt]
    Portfolio: \href{https://www.nabilahmad.site}{www.nabilahmad.site} \quad$|$\quad GitHub: \href{https://github.com/nabilroghani}{github.com/nabilroghani}
\end{center}
\vspace{-4pt}

%----------------------------------------------------------------------------------------
%   SECTION 1
%----------------------------------------------------------------------------------------
\section{1. Academic Background \& Motivation for European Graduate Research}
I am writing to express my strong candidacy for admission to the Master of Science (MSc) Program in Computer Science at \textbf{[UNIVERSITY\_NAME]}. As a graduating Computer Science student from the University of Peshawar (CGPA: \textbf{3.32 / 4.00}, European ECTS First-Class equivalent; 100\% English-medium instruction), my academic trajectory has been defined by a commitment to bridging software engineering rigor with intelligent, human-centered systems. Europe’s progressive leadership in trustworthy AI, European Union digital accessibility mandates, and open scientific inquiry provides the premier intellectual landscape to advance my research in \textbf{Human-Centered Software Engineering, Accessibility, and Applied AI}.

%----------------------------------------------------------------------------------------
%   SECTION 2
%----------------------------------------------------------------------------------------
\section{2. International Industry Experience with European Clients \& Standards}
Unlike traditional candidates with purely classroom foundations, I bring substantive remote engineering experience delivering commercial software systems for \textbf{European clients, particularly in Malta and the United Kingdom}. At Apptex Software Solutions, I architected production applications requiring adherence to European regulatory and quality benchmarks. On \textbf{Accurit Accountancy} (a Malta-based firm), I implemented accessibility engineering compliant with \textbf{WCAG 2.2 Level AA} guidelines, using automated axe-core accessibility auditing tools alongside comprehensive end-to-end (E2E) testing suites in Playwright (Chromium, Firefox, WebKit) and Vitest. This experience instilled a deep respect for digital inclusion, robust test coverage, and fault-tolerant architecture.

My international portfolio further spans systems like \textbf{Galea's Kitchen} (Malta hospitality, live CeleryHQ reservation API integration), \textbf{Verbalisers} (European translation services quote platform), and \textbf{Baratie Stays} (Malta vacation rentals). Concurrently, I explored applied AI by architecting \textbf{MedPulse AI}, an experimental healthcare triage system coupling Google Gemini 2.0 with a MongoDB-grounded Retrieval-Augmented Generation (RAG) pipeline to generate automated bilingual (English/Urdu) medical explanations. These experiences demonstrated that cutting-edge AI requires strict software engineering validation to be viable in critical human environments.

%----------------------------------------------------------------------------------------
%   SECTION 3
%----------------------------------------------------------------------------------------
\section{3. Research Objectives \& Alignment with [RESEARCH\_GROUP]}
I am intensely motivated to join the \textbf{[RESEARCH\_GROUP]} under the supervision of Prof.~\textbf{[PROFESSOR\_NAME]} at \textbf{[UNIVERSITY\_NAME]}. My proposed research focuses on \textbf{Accessible, Trustworthy Human-AI Interfaces and Automated Software Verification}. As generative AI models become embedded in public and commercial platforms, ensuring they comply with accessibility standards (such as the European Accessibility Act) and provide verifiable, transparent responses is paramount. My technical background in modern web architectures (Next.js 16, TypeScript), automated testing frameworks (Playwright, Vitest), and vector-grounded RAG pipelines enables me to immediately assist ongoing laboratory investigations, rapidly prototyping interactive testbeds and designing empirical experiments with end users.

%----------------------------------------------------------------------------------------
%   SECTION 4
%----------------------------------------------------------------------------------------
\section{4. Long-Term Career Vision in Europe}
During my Master’s studies, I intend to publish my research at leading international conferences in software engineering and HCI. Following my MSc, I aim to pursue a Ph.D. or join a research-focused European industrial laboratory, contributing to the development of ethically grounded, highly reliable intelligent software systems. Having collaborated remotely with European clients for years, I am well-prepared to integrate seamlessly into a multicultural European university environment, bringing relentless dedication, technical independence, and collaborative energy to your research group.

\end{document}

